\section{Trabalhos Relacionados}

O problema tratado neste artigo dialoga com quatro frentes. A primeira é a infraestrutura legislativa aberta no Brasil, especialmente o SAPL e o ecossistema do Interlegis, que padronizam parte da tramitação e da publicidade legislativa \cite{interlegis_sapl}. Essa camada institucional é relevante porque viabiliza coleta automatizada, mas não resolve, por si só, a heterogeneidade de implantação observada no nível municipal.

A segunda frente é a recuperação de informação em documentos legislativos e jurídicos brasileiros. No nível federal, já há trabalhos que estudam stemming, BM25, pipelines de recuperação e \textit{relevance feedback} para a Câmara dos Deputados \cite{souza_stemming_2021,vitorio_ir_pipeline_2021,vitorio_rfcorpus_2025}. Em paralelo, recursos como UlyssesNER-Br mostram amadurecimento do ecossistema brasileiro de NLP jurídico, com novos corpora e tarefas anotadas em português \cite{albuquerque_ulyssesner_2022}.

A terceira frente é a literatura de difusão e transferência de políticas públicas, segundo a qual governos aprendem, imitam ou adaptam soluções observadas em outras jurisdições \cite{shipan_volden_2008}. Embora essa literatura não trate de processamento de linguagem natural, ela fundamenta conceitualmente a utilidade de comparar proposições entre cidades e de transformar acervos dispersos em infraestrutura de aprendizado intermunicipal.

Essas linhas são próximas ao problema aqui tratado, mas ainda insuficientes para cobri-lo integralmente. Em geral, elas assumem um acervo previamente disponível, mais homogêneo e frequentemente concentrado no nível federal ou em domínios jurídicos já curados. No cenário municipal, o desafio começa antes da busca: é necessário descobrir as fontes, lidar com a heterogeneidade dos portais, extrair os dados em escala e normalizar textos legislativos curtos e ruidosos antes de qualquer comparação.

Embora existam contribuições relevantes sobre infraestrutura legislativa, recuperação de informação jurídica e difusão de políticas, ainda falta uma abordagem que integre descoberta automatizada de fontes, extração em escala, textualização e recuperação semântica no contexto municipal brasileiro. Em relação aos trabalhos revisados, a principal contribuição do Sonar Municipal é integrar descoberta de fontes, extração, textualização, busca semântica e associação exploratória com indicadores em um único pipeline orientado ao nível municipal.
